\documentclass{article}
\usepackage{kotex}
\usepackage{setspace}  % 줄간격
\setstretch{1.15}
\usepackage{enumitem}  % 목록 들여쓰기 
\setlist[itemize]{leftmargin=*, align=left}
\setlist[description]{
  font=\normalfont,
  labelwidth=4cm,
  leftmargin=4cm,
  labelsep=0cm,
  align=left
}
\usepackage{titlesec}  % 섹션 스타일 조정
\titleformat{\section}
  {\normalfont\Large\bfseries}
  {}
  {0pt}
  {}
  [\vspace{0.3em}\titlerule]

\titlespacing*{\section}
  {0pt}
  {1.5em}
  {0.8em}
\newcommand{\descriptionrule}{%
  \item[]\vspace{-0.6em}
  \noindent\hspace*{-\leftmargin}%
  \rule{0.6\linewidth}{0.4pt}
  \vspace{-0.4em}
}              
  \usepackage[a4paper, margin=2.5cm]{geometry}
\usepackage{fancyhdr}                 
\pagestyle{fancy}
\fancyhf{}                             
\fancyfoot[L]{\nouppercase{음악과 사회 (2026학년도 2학기)}} 
\fancyfoot[R]{\thepage}                
\renewcommand{\headrulewidth}{0pt}    
\renewcommand{\footrulewidth}{0.4pt}   

\begin{document}

\title{음악과 사회 (2026학년도 2학기)}
\author{\textit{written by using} \LaTeX}
\date{\today}
\maketitle
\thispagestyle{fancy}
\setcounter{secnumdepth}{0}
\bigskip

\section{0. 강좌정보}
\medskip
교과목번호: C10.146(002)\\
개설학과: 서울대학교 음악대학 음악학과\\
담당교수: 이서현\\
수업시간: 월/수 16:00--17:15\\
수업장소: 220-203

\section{1. 수업목표}
\medskip
음악의 사회적, 문화적, 상업적, 경제적 측면에 초점을 맞춰 음악의 내용을 체계적으로 학습한다. 예술 음악의 장르와 함께, 재즈, 블루스, 컨트리, 뮤지컬, 록, 팝, 포크, 디스코, 소울, 힙합, 트로트 등의 다양한 대중음악 장르를 다룬다. 각 장르의 역사적 사회적 배경에 대한 지식과 더불어, 실제로 음악 레퍼토리를 분석함으로써, 음악의 형식, 화성, 음계, 리듬, 연주 등의 음악적 내용 뿐 아니라 음악의 제작, 생산, 유통, 소비 등의 사회적 경제적 이해를 도모한다.

\section{2. 교재 및 참고문헌}
\medskip
N/A

\newpage
\section{3. 주차별 강의계획}
\medskip
\begin{description}
    \item[1주차] 강의 소개 및 오리엔테이션
    \item[2주차] 대중음악: 팝(pop)
    \item[3주차] 컨트리(country)
    \item[4주차] 블루스(blues)
    \item[5주차] 재즈(jazz)
    \item[6주차] 뮤지컬(musical)
    \item[7주차] 흑인 음악
    \item[8주차] 자유학습일 - 중간 과제 작성
    \item[9주차] 록(rock)
    \item[10주차] 포크 (folk)
    \item[11주차] 댄스음악
    \item[12주차] 월드 뮤직 I : 라틴 아메리카
    \item[13주차] 월드뮤직 II: 유럽
    \item[14주차] 월드 뮤직 III: 아프리카
    \item[15주차] 기말고사
\end{description}

\section{4. 평가방법}
\medskip
\begin{description}
    \item[\textbf{성적부여방식}] 상대평가(A--F)
    \item[\textbf{출석(10\%)}] 수업일수의 1/3을 초과하여 결석하면 성적은 ``F'' 또는 ``U''가 됨\\(담당교수가 불가피한 결석으로 인정하는 경우는 예외로 할 수 있음)
    \item[\textbf{과제(20\%)}] 주별 감상 과제
    \item[\textbf{중간(15\%)}] 내 인생에서 의미가 있었던 음악들을 선곡하고 설명(분량 자유, ~11/4(수))
    \item[\textbf{기말(40\%)}] 12/14 시행
    \item[\textbf{수시평가(15\%)}]  
    \descriptionrule
    \item[\textbf{합계}] 100\%
\end{description}

\newpage
\section{5. 생성형 AI 도구 활용 방침}
\medskip
본 강의에서는 생성형 AI를 학습 보조 도구로 활용할 수 있습니다. 과제 수행에 필요한 아이디어 및 주제 탐색, 자료분석 등의 목적으로 생성형 AI를 사용할 수 있습니다. 다만, 생성형 AI를 과제 수행에 사용한 경우, 사용 도구명, 사용 과정 및 생성 내용 등을 과제에 명시해야 합니다.

\section{6. 정원 외 신청 수용 여부}
\medskip
최대 0명

\section{7. 장애학생 지원사항}
\medskip
\subsection{강의 수강 관련}
\medskip
\begin{itemize}
    \item 시각장애: 교재 제작(디지털교재, 점자교재, 확대교재 등), 대필도우미 허용
    \item 지체장애: 교재 제작(디지털교재), 대필도우미 및 수업보조 도우미 허용
    \item 청각장애: 대필 및 문자통역 도우미 활동 허용, 강의 녹취 허용
    \item 건강장애: 질병 등으로 인한 결석에 대한 출석 인정, 대필도우미 허용
    \item 학습장애: 대필도우미 허용
    \item 지적장애/자폐성장애: 대필도우미 및 수업 멘토 허용
\end{itemize}
\subsection{과제 및 평가 관련}
\medskip
\begin{itemize}
    \item 시각장애/지체장애/청각장애/건강장애/학습장애: 과제 제출기한 연장, 과제 제출 및 응답 방식의 조정, 평가 시간 연장, 평가 문항 제시 및 응답 방식의 조정, 별도 고사실 제공
    \item 지적장애/자폐성장애: 개별화 과제 제출 및 대체 평가 실시
\end{itemize}
\subsection{비고}
\medskip
본 강의를 수강하는 장애학생들에게는 이상의 지원 서비스 이외에도 장애학생 개개인의 특성과 요구에 따라, 지도교수 및 장애학생지원센터와의 상담을 통하여 적절한 수준의 지원 서비스를 제공합니다. 장애학생에 대한 지원서비스와 관련하여 문의사항이 있는 학생들은 담당교수 이서현 혹은 장애학생지원센터(02-880-8787)로 문의바랍니다.

\end{document}
